% TOP_PROBLEM: {"id": 30006746, "title": "No-Slip Vanishing-Viscosity Limit in the Energy Norm", "category_id": 9, "category": {"id": 9, "name": "pde", "display_name": "Partial Differential Equations", "description": "PDEs and their applications in physics and geometry.", "slug": "pde", "order_index": 9, "created_at": "2026-07-31T15:26:25.671Z"}, "status": "open"}
% FIRST_POSTED: null
% !TeX program = lualatex
% Compile twice: lualatex 30006746.tex
% All references and prose are embedded; no BibTeX database or external inputs.
\documentclass[11pt,a4paper]{article}
\usepackage[margin=24mm,headheight=14pt]{geometry}
\usepackage{fontspec}
\setmainfont{Latin Modern Roman}
\setsansfont{Latin Modern Sans}
\setmonofont{Latin Modern Mono}
\usepackage{amsmath,amssymb,mathtools}
\usepackage{unicode-math}
\setmathfont{Latin Modern Math}
\usepackage{microtype}
\usepackage{enumitem,xurl,needspace}
\usepackage[dvipsnames]{xcolor}
\usepackage{hyperref}
\hypersetup{unicode=true,colorlinks=true,linkcolor=MidnightBlue,urlcolor=MidnightBlue,
 pdftitle={Research attempt: Problem 30006746},pdfauthor={Research notebook},bookmarksnumbered=true}
\usepackage{bookmark}
\usepackage{tocloft}
\setlength{\cftsecnumwidth}{3em}
\cftsetpnumwidth{2.5em}
\cftsetrmarg{4em}
\usepackage{titlesec}
\titleformat{\section}{\Large\bfseries\raggedright}{\thesection}{0.7em}{}
\usepackage{fancyhdr}
\pagestyle{fancy}\fancyhf{}
\fancyhead[L]{\small\sffamily Open Problems Math | Research notebook}
\fancyhead[R]{\small Problem 30006746 | 27 September 2026}
\fancyfoot[C]{\thepage}
\setlength{\parindent}{0pt}
\setlength{\parskip}{3pt plus 1pt}
\setlength{\emergencystretch}{3em}
\setcounter{tocdepth}{1}
\setcounter{secnumdepth}{1}
\setlist[itemize]{leftmargin=3em,itemsep=5pt,topsep=4pt,parsep=0pt}
\allowdisplaybreaks
\newcommand{\N}{\mathbb{N}}
\newcommand{\Z}{\mathbb{Z}}
\newcommand{\Q}{\mathbb{Q}}
\newcommand{\R}{\mathbb{R}}
\newcommand{\C}{\mathbb{C}}
\newcommand{\F}{\mathbb{F}}
\newcommand{\A}{\mathbb{A}}
\newcommand{\PP}{\mathbb P}
\newcommand{\E}{\mathbb{E}}
\newcommand{\eps}{\varepsilon}
\newcommand{\dd}{\,\mathrm d}
\newcommand{\sref}[2]{\hyperlink{source:#1:#2}{[S#2]}}
\newcommand{\eref}[2]{\hyperlink{extra:#1:#2}{[E#2]}}
% Turn the next switch off for a shorter reading copy. The quotations preserve
% every exactTarget field, including qualifications omitted from a short summary.
\newif\ifshowcatalogue
\showcataloguetrue
\newcommand{\cataloguescope}[1]{\ifshowcatalogue
 \paragraph{Exact scope in the supplied catalogue.}
 \begin{quote}\small #1\end{quote}\fi}
\usepackage{etoolbox}
\AtBeginEnvironment{itemize}{\small}
\numberwithin{equation}{section}
% Standalone export. Original research content preserved.
\begin{document}
\setcounter{section}{0}
% ================================================================
% problemId: 30006746
% Canonical problem ID: 30006746
% exactTarget SHA-256: 7f4370827d396cd852db19e6604f2ef6fbf0bc778845a201b55509c122696e32
\clearpage
\section*{\texorpdfstring{No-Slip Vanishing-Viscosity Limit in the Energy Norm}{No-Slip Vanishing-Viscosity Limit in the Energy Norm}}\label{30006746}
\subsection{Definitions and mathematical statement}
For fixed smooth divergence-free $u_0$ vanishing on $\partial\Omega$, let $u^\nu$ be any no-slip Leray--Hopf weak solution of Navier--Stokes and let $u$ be the smooth Euler solution with $u\cdot n=0$. A Leray--Hopf solution has finite-energy weak evolution and satisfies the energy inequality. For each interval $[0,T]$ of smooth Euler existence, the target is
\[
 \lim_{\nu\downarrow0}\operatorname*{ess\,sup}_{0\le t\le T}
 \|u^\nu(t)-u(t)\|_{L^2(\Omega)}=0.
\]
Both $d=2$ and $d=3$ are included. The same initial data are used at every viscosity, and convergence is asserted for every allowed weak-solution family, not a specially chosen subsequence.
\par\smallskip\textit{Formulation and concept references:} \sref{30006746}{1}.
\cataloguescope{For d in \{2,3\}, every smooth bounded domain Omega in R\textasciicircum{}d and every fixed smooth divergence-free initial velocity u0 vanishing on the boundary, prove or disprove that every no-slip Leray–Hopf Navier–Stokes family u\textasciicircum{}nu with viscosity nu and initial value u0 converges to the Euler solution u in L\textasciicircum{}infinity([0,T];L\textasciicircum{}2(Omega)) as nu tends to zero, for each T on which u is smooth and satisfies u.n=0.}
\subsection{Short English statement}
Do all no-slip viscous flows converge in the uniform-in-time energy norm to the smooth inviscid flow as viscosity vanishes?
% BEGIN RESEARCH ADDITION 136
\subsection{Research attempt}

\paragraph{Current frontier.}
Kato's criterion identifies the missing boundary estimate: in the fixed-data,
smooth-Euler setting, the requested convergence is equivalent to
\begin{equation}\label{30006746:kato}
 D_\nu(c):=\nu\int_0^T\int_{d(x,\partial\Omega)<c\nu}
                       |\nabla u^\nu|^2\,dx\,dt\longrightarrow0
\end{equation}
for fixed $c>0$. It is an equivalence applicable to energy-inequality
solutions, not a theorem that this expression always vanishes. The supplied
survey states it as Theorem 3.7. Instability constructions using
viscosity-dependent initial boundary layers do not automatically meet this
question's single fixed datum; see the discussion after its Theorem 3.6.
\sref{30006746}{1}

\paragraph{Attack route.}
Testing the viscous equation against Euler works if its full boundary trace
remains zero. A boundary corrector works under \eqref{30006746:kato}. Stronger
integrability would force that condition if its viscosity loss is below an
explicit threshold. We investigate the first and third routes, and test the
proposed shortcut that initially no-slip Euler flow stays no-slip.

\paragraph{Attempt: the trace-preserving special case.}
Assume, in addition to the actual hypotheses, that
$u|_{\partial\Omega}=0$ for the whole smooth interval. A Leray--Hopf solution
with the same initial value can then be tested weakly against $u$.
Combining that identity, Euler energy conservation and the energy inequality
for $u^\nu$ gives, for $w=u^\nu-u$,
\[
 \frac12\|w(t)\|_2^2+\nu\int_0^t\|\nabla u^\nu\|_2^2
 \le\int_0^t\|\nabla u\|_\infty\|w\|_2^2
   +\nu\int_0^t\langle\nabla u^\nu,\nabla u\rangle.
\]
Time approximation justifies the smooth time-dependent test; no uniqueness
of $u^\nu$ is used. Young's inequality and Gronwall give
\begin{equation}\label{30006746:trace}
 \operatorname*{ess\,sup}_{t\le T}\|u^\nu(t)-u(t)\|_2^2
 \le\nu\exp\!\left(2\int_0^T\|\nabla u\|_\infty dt\right)
                  \int_0^T\|\nabla u\|_2^2dt.
\end{equation}
This proves a uniform $O(\sqrt\nu)$ energy-norm convergence estimate for
this restricted class and every one of its allowed weak-solution families.
It is a direct relative-energy calculation, not a claim of novelty.

\paragraph{A fixed smooth datum that immediately loses the full Euler trace.}
In the unit disk set $s=x^2+y^2$, $\psi=x(1-s)^2$, and
$u_0=(-\partial_y\psi,\partial_x\psi)$. Both components vanish on $s=1$,
and $\nabla\cdot u_0=0$. The initial pressure solves
\[
 \Delta p_0=-\operatorname{tr}((Du_0)^2),\qquad \partial_n p_0=0
 \quad(s=1).
\]
An explicit polynomial solution, up to an irrelevant constant, is
\begin{align*}
 p_0={}&-\frac35s^4+\frac73s^3-\frac{10}3s^2+\frac{31}{15}s\\
 &+x^2\left(\frac{16}5s^3-10s^2+\frac{32}3s-\frac{62}{15}\right).
\end{align*}
Differentiation verifies the equations exactly; the companion script checks
them in $\Q[x,y]$. On the circle,
\[
 p_0=\frac15+\frac4{15}y^2,\qquad
 (-y\partial_x+x\partial_y)p_0=\frac8{15}xy.
\]
Since $(u_0\cdot\nabla)u_0=0$ at the boundary, Euler's equation gives
\[
 \partial_tu(0)\cdot(-y,x)=-\frac8{15}xy,
\]
which is not identically zero. Local smooth Euler flow from this polynomial
datum therefore immediately loses its tangential zero trace, while keeping
its normal trace zero. This is not an inviscid-limit counterexample; it
refutes the extra trace-persistence lemma needed to apply
\eqref{30006746:trace} universally.

\paragraph{A quantitative boundary-integrability reduction.}
For $p>2$, the smooth tubular-neighborhood volume estimate and spatial
H\"older inequality give
\[
 D_\nu(c)\le C_{\Omega,c,p}\nu^{2-2/p}
                    \int_0^T\|\nabla u^\nu(t)\|_p^2dt.
\]
Hence the additional, uniform estimate
\begin{equation}\label{30006746:missing}
 \int_0^T\|\nabla u^\nu\|_p^2dt=O(\nu^{-\alpha}),
 \qquad \alpha<2-2/p,
\end{equation}
for every allowed family would prove the exact target by Kato's criterion.
The strip $L^p$ norm would suffice. At $p=2$, the energy inequality gives
$O(\nu^{-1})$, precisely the nondecaying threshold, not a strict improvement.
Thus the available energy estimate cannot be substituted into this step.

\paragraph{Stress test.}
A boundary layer with order-one velocity jump across thickness $\nu$ has
squared $L^2$ mass of order $\nu$ but order-one contribution to
\eqref{30006746:kato}. This is a dimensional test, not a Navier--Stokes solution.
Small layer mass therefore cannot by itself exclude boundary dissipation.
The pressure calculation concerns actual fixed smooth data but says nothing
about how much its viscous boundary layer dissipates. No unproved uniqueness
or compactness of arbitrary three-dimensional weak solutions was used in
\eqref{30006746:trace}. The $p$-integrability in \eqref{30006746:missing} is an extra
hypothesis, not a property deduced from the Leray--Hopf definition.

\paragraph{Outcome.}
\textbf{Outcome: Conditional reduction.}
The attempt gives a sufficient integrability threshold, a relative-energy
proof when Euler retains no-slip, and an exact polynomial countertest to
automatic trace persistence. For the full problem, one still needs Kato
boundary dissipation, an estimate such as \eqref{30006746:missing}, or a genuine
fixed-data family violating convergence. None is established here.

% END RESEARCH ADDITION 136
\subsection{Sources}
\begin{itemize}\raggedright
\item[\textnormal{[S1]}]\hypertarget{source:30006746:1}{} Y. Maekawa and A. Mazzucato, The Inviscid Limit and Boundary Layers for Navier-Stokes Flows (2016).
\url{https://arxiv.org/html/1610.05372}.
{\small\textit{Catalogue formulation source.}}
% BEGIN RESEARCH REFERENCES 136
% The supplied survey, particularly Theorems 3.6--3.7, is the source used.

% END RESEARCH REFERENCES 136
\end{itemize}

\end{document}
